\chapter{\texorpdfstring{The Unified Equation: Derivazione di
\(T_s^{(2)}\)}{The Unified Equation: Derivazione di T\_s\^{}\{(2)\}}}\label{lequation-unificata-followszione-di-t_s2}

\begin{quote}
\textbf{Core argument:} la stabilità di un sistema egemonico è
measurable come rapporto strutturato tra forze di coerenza e forze di
decoherence, modulato da due variables --- antifragility (\(A\)) e
involution (\(\Omega\)) --- che agiscono su dimensioni diverse del
sistema.
\end{quote}



\section{Premessa di Coerenza
Notazionale}\label{premessa-di-coerenza-notezionale}

Prima di followsre l'equation, fissiamo in modo univoco la notezione ---
che nel corpus DEQ precedente oscillava tra due convenzioni.

\Hypoth{} \textbf{Convenzione unificata adottata in questo volume:}

Il \textbf{discount factor} \(\delta \in [0, 1]\) è un \emph{orizzonte
temporale}: misura quanto il sistema valuta il futuro with respect tol
presente. È \textbf{stabilizzante}. Nel Folk Theorem della
\emph{Geometria dell'Egemonia}, la cooperazione sostenibile richiede
\(\delta \geq \delta^* = 0{,}50\). Alto \(\delta\) = bassa preferenza
temporale = cooperazione possibile = stabilità.

\Proved{} \textbf{Confollowsnza formal:} \(\delta\) appare \emph{al
numeratore} del \(T_s\). Un sistema con \(\delta\) alto è più coerente.

L'incoerenza del \emph{Memoriale Tecnico DEQ} (2026-04-23), dove
\(\delta\) compariva al denominatore, è qui risolta: quella era
un'abbreviazione tipografica della quantità complementare
\((1-\delta)\), che misura la \textbf{preferenza per il presente}
(destabilizzante). Per chiarezza assoluta, in questo volume usiamo
sempre \(\delta\) come orizzonte stabilizzante.



\section{La Doppia Eredità: Geometria e
Topografia}\label{la-doppia-eredituxe0-geometria-e-topografia}

\subsection{Il Payoff Macro della Geometria
dell'Egemonia}\label{il-payoff-macro-della-geometria-dellegemonia}

Nel volume \emph{La Geometria dell'Egemonia} (Marcelino, 2026), l'attore
\(i\) ottiene payoff di periodo:

\[
\Pi_i(t+1) = H_i \cdot Q_i \cdot (1 - \delta_i \cdot s_{i,3})
\]

where \(H_i\) è l'impronta sistemica,
\(Q_i = (s_{i,1}+s_{i,2}+s_{i,3})/\sqrt{3}\) la qualità strategica, e
\(s_{i,3}\) la risorsa relazionale (legitimacy).

\subsection{La Funzione Predittiva Micro della Topografia del
Discorso}\label{la-funzione-predittiva-micro-della-topografia-del-discorso}

Nel volume \emph{La Topografia del Discorso Emancipativo} (Marcelino,
2026, \S{}5.7), per un singolo coupling emittente-ricevente:

\[
\Pi(t+1) = M_e \cdot Q_0 \cdot (1 - \delta_R \cdot z_R)
\]

where \(M_e\) è la communicative mass dell'emittente,
\(Q_0 = (x+y+z)/\sqrt{3}\) la ethical quality del discorso, \(\delta_R\)
il discount factor del ricevente, \(z_R\) la sua complessità
dialettica.

\subsection{\texorpdfstring{\Proved{} L'Identità Formale (Risultato
Strutturale)}{ L'Identità Formale (Risultato Strutturale)}}\label{lidentituxe0-formal-risultato-structural}

Le due formule \textbf{non sono analoghe --- sono la stessa equation
applicata a due scale diverse}:

{\def\LTcaptype{none} % do not increment counter
{\small\begin{longtable}[]{@{}
  >{\raggedright\arraybackslash}p{(\linewidth - 4\tabcolsep) * \real{0.3333}}
  >{\raggedright\arraybackslash}p{(\linewidth - 4\tabcolsep) * \real{0.3333}}
  >{\raggedright\arraybackslash}p{(\linewidth - 4\tabcolsep) * \real{0.3333}}@{}}
\toprule\noalign{}
\begin{minipage}[b]{\linewidth}\raggedright
Grandezza
\end{minipage} & \begin{minipage}[b]{\linewidth}\raggedright
Topografia (micro-comunicativo)
\end{minipage} & \begin{minipage}[b]{\linewidth}\raggedright
Geometria (macro-egemonico)
\end{minipage} \\
\midrule\noalign{}
\endhead
\bottomrule\noalign{}
\endlastfoot
Massa di proiezione & \(M_e\) & \(H_i\) \\
Qualità sull'asse etico & \(Q_0 = (x+y+z)/\sqrt{3}\) &
\(Q_i = (s_{i,1}+s_{i,2}+s_{i,3})/\sqrt{3}\) \\
Fattore di erosione & \((1 - \delta_R z_R)\) &
\((1 - \delta_i s_{i,3})\) \\
Soglia di cooperazione (Folk Theorem) & \(\delta^*_E = 0{,}50\) &
\(\delta^* = 0{,}50\) \\
\end{longtable}}
}

\Proved{} \textbf{Confollowsnza:} la coincidenza numerica
\(\delta^* = 0{,}50\) tra Discursive Folk Theorem (Topografia App.~A) e
Geometry Folk Theorem \emph{non è una coincidenza} --- è
l'scale invariance di un'unica struttura matematica. Questo è il
primo risultato structural della DEQ$^2$: la stessa legge governa il
pagamento dell'emittente verso il ricevente e il pagamento dell'egemone
verso l'attore subordinato.

\Hypoth{} \textbf{Mossa di passaggio al livello sistemico:} non
sostituiamo la logica dell'attore singolo --- la \emph{integriamo}. Il
payoff individuale \(\Pi_i(t+1)\) resta valido come legge local; ciò
che cerchiamo è la quantità che descrive lo \emph{stato collettivo},
\(T_s\), come ordine emergesnte dal campo \(\{\Pi_j\}\).



\section{3.1bis Il Bridge Micro/Macro via
Kuramoto}\label{bis-il-bridge-micromacro-via-kuramoto}

\Hypoth{} Considereremo il sistema come una popolazione di \(N\)
accoppiamenti elementari (cittadino-discorso, attore-egemone,
élite-massa) ciascuno con instantaneous payoff \(\Pi_j(t)\) e fase
\(\theta_j(t)\) (la posizione nel ciclo di legittimazione/decoherence).
Definiamo il \textbf{order parameter di Kuramoto} (1984):

\[
\Phi(t) e^{i\psi(t)} = \frac{1}{N}\sum_{j=1}^{N} e^{i\theta_j(t)} \qquad \Phi \in [0, 1]
\]

\begin{itemize}
\tightlist
\item
  \(\Phi \to 1\): tutte le fasi allineate (coerenza piena);
\item
  \(\Phi \to 0\): fasi distribuite uniformemente (decoherence piena).
\end{itemize}

\Hypoth{} \textbf{Operatore di trasposizione:} definiamo la
\emph{quantità sistemica} followsta da una local quantity \(X_j\) come

\[
\langle X \rangle_\Phi := \Phi \cdot \frac{1}{N}\sum_j X_j
\]

L'aggregazione ``naive'' \(\frac{1}{N}\sum X_j\) è recuperata per
\(\Phi = 1\); in regime decoerente \(\Phi < 1\) il sistema \emph{perde
sostanza} anche a parità di payoff individuali. È questa la differenza
tra un USA decoerente con alti payoff individuali e un Brazil
\(\Phi\)-engine con payoff individuali più modesti ma allineati.

\Proved{} \textbf{Risultato:} \(T_s^{(2)}\) followsto nelle sezioni
successive è formalmente \(\langle T_s^{\text{loc}}\rangle_\Phi\) del
campo Topografia, modulato da \(A\) e \(\Omega\). La forma sistemica
\(T_s^{(2),\text{sys}}\) del Ch.~4 esplicita questo passaggio.



\section{\texorpdfstring{The Evolutionary Foundation of \(\Phi\): Captured Cooperation}{The Evolutionary Foundation of Phi: Captured Cooperation}}\label{il-fondamento-evolutivo-di-phi}

\Hypoth{} \textbf{The causal problem.} Section~3.1bis introduced
\(\Phi\) as the Kuramoto order parameter --- a mathematical quantity
measuring phase alignment. The causal question remains open:
\emph{why} do phases align or decohere? What is the mechanism that
produces \(\Phi > 0\) as a stable state, and decoherence as its
engineered counterpart?

The answer comes from evolutionary biology: cooperation is the natural
state when the mechanisms sustaining it outweigh those eroding it.
\(\Phi \to 0\) is not the baseline equilibrium --- it is a state
\emph{engineered} by systematic extraction.

\Proved{} \textbf{Hamilton's condition (1964).} Cooperation is
evolutionarily stable in a population if and only if:

\[
rB > C \tag{H}
\]

where \(r\) is the relatedness coefficient, \(B\) is the benefit to
the recipient, and \(C\) is the cost to the donor. This condition ---
derived for biology and applied here as a structural isomorphism
(\Hypoth{}) --- has its geopolitical analogue: systemic cooperation is
stable when collective benefit exceeds the individual cost of cooperating,
weighted by the structural coherence of the system.

\Hypoth{} \textbf{Structural correspondence with \(\Phi\).}
The order parameter \(\Phi > 0\) is the systemic signature of the
\(rB > C\) regime: phases align when cooperative mechanisms
(reciprocity, institutional norms, trust) outweigh defection mechanisms.
This is not an identity --- it is a regime correspondence:

\[
rB > C \;\xleftrightarrow{\;\text{\Hypoth{}}\;}\; \Phi > 0
\]

\Hypoth{} \textbf{The CPGG as capture mechanism.} The
\textit{Captured Public Goods Game} (CPGG --- §D.4) formalizes the
mechanism by which cooperation is eroded by an extractive sub-network
of fraction \(\phi \ll 1\) with endogenous positive feedback
(capture loop, Proposition~1 \Proved{}):

\[
|\mathcal{E}| \uparrow \;\Rightarrow\; \sigma_{\text{san}} \downarrow \;\Rightarrow\; \pi_E \uparrow \;\Rightarrow\; |\mathcal{E}| \uparrow
\]

The effective degree of cooperators on scale-free networks (the topology
of real-world networks) collapses with \(\phi\):

\[
k_{\text{eff}}(\phi) = \bar{k} \cdot \frac{1 - \phi^{(\gamma-2)/(\gamma-1)}}{1-\phi} \quad (\text{\Proved{} F3})
\]

The critical threshold \(\phi^*\) is the point at which the condition
\(b/c > k_{\text{eff}}(\phi)\) ceases to hold and \(\Phi\) collapses.
The system signals the transition \emph{before} cooperative collapse:
monitoring \(\phi\) provides an early warning indicator of
\(T_s^{(2),\text{sys}} \to 0\).

\Proved{} \textbf{The multilevel selection foundation.} The causal
justification for the micro/macro level shift is provided by
\textit{multilevel selection theory} (Wilson \& Sober 1994; Okasha 2006).
The Price equation (1970) --- an exact algebraic identity --- decomposes
the variation in cooperative behaviour into an inter-group term
(coherence) and an intra-group term (extraction):

\[
\Delta\bar{z} = \underbrace{\mathrm{Cov}(W_k,\bar{z}_k)}_{\text{inter-group}} + \underbrace{\mathbb{E}_k\!\left[\mathrm{Cov}(W_{ik},z_{ik})\right]}_{\text{intra-group}} \tag{Price 1970}
\]

\Hypoth{} The structural correspondence is: the inter-group term is
proportional to \(\Phi > 0\); the intra-group term is proportional to
\(-\Omega\).

This structure justifies the form of \(T_s^{(2),\text{sys}}\): the
numerator amplifies inter-group selection (coherence, associated with
\(\Phi\langle A\rangle\)), the denominator amplifies intra-group
selection (extraction, associated with \(\langle\Omega\rangle/\Phi\)).

\textit{For the complete derivation of the evolutionary framework, see §D.4.}



\section{\texorpdfstring{La Soglia di Sgretolamento \(T_s\)
(DEQ$^1$)}{La Soglia di Sgretolamento T\_s (DEQ$^1$)}}\label{la-soglia-di-sgretolamento-t_s-dequxb9}

La first formulazione della DEQ, presentata nel \emph{Memoriale Tecnico}
(2026-04-23), pone:

\[
T_s^{(1)} = \frac{s_{iii} \cdot z \cdot \delta}{\sigma \cdot \varepsilon_{\text{dis}}}
\]

where:

\begin{itemize}
\tightlist
\item
  \(s_{iii}\) è la \textbf{ethical-discursive legitimacy} del sistema
  (ereditata dalla Geometria);
\item
  \(z\) è la \textbf{projection on the dialectical axis} (complessità del
  discorso pubblico --- eredità della Topografia);
\item
  \(\delta\) è l'\textbf{time horizon} (Folk Theorem ---
  Geometria \S{}3);
\item
  \(\sigma\) è la \textbf{stochastic volatility} ereditata
  dall'estensione EDSD (Geometria \S{}12-13);
\item
  \(\varepsilon_{\text{dis}} \in [0,1]\) è il \textbf{coefficiente di
  disimpegno morale}, ispirato a Bandura (1999) e Milgram (1974).
\end{itemize}

\Proved{} \textbf{Lettura formal:} al numeratore stanno le forze di
coerenza (legitimacy $\times$ complessità $\times$ orizzonte); al denominatore le
forze di decoherence (volatilità $\times$ disimpegno). \(T_s^{(1)} = 1\) è la
critical threshold. \(T_s^{(1)} \gg 1\) indica sistema stabile;
\(T_s^{(1)} < 1\) sistema in phase transition imminente.

\Conj{} \textbf{Limite della DEQ$^1$:} l'equation \emph{non distingue} tra
sistemi semplicemente resilienti e sistemi antifragili, né tra sistemi
sotto attrito esogeno (volatilità esterna) e sistemi sotto attrito
endogeno (competizione involutiva). Due sistemi con lo stesso
\(T_s^{(1)}\) possono avere traiettorie radicalmente diverse. Questo
motiva l'estensione DEQ$^2$.



\section{I Due Assi Strutturali
Mancanti}\label{i-due-assi-strutturali-mancanti}

Introduciamo due variables che catturano dimensioni non incluse nella
DEQ$^1$.

\subsection{\texorpdfstring{L'Antifragility
\(A\)}{L'Antifragility A}}\label{lantifragilituxe0-a}

\Proved{} \textbf{Definizione formal (da Taleb 2012, formalizzazione
structural):}

\[
A = \left.\frac{\partial^2 \Pi}{\partial \sigma^2}\right|_{\sigma_0}
\]

where \(\Pi\) è il payoff di sistema e \(\sigma_0\) è il livello base di
volatilità. \(A\) misura la \textbf{convessità della risposta allo
stress}:

\begin{itemize}
\tightlist
\item
  \(A > 0\): positive convexity $\Rightarrow$ il sistema \emph{guadagna} dalla
  volatilità (antifragile);
\item
  \(A = 0\): linearità $\Rightarrow$ sistema resilient neutro;
\item
  \(A < 0\): concavità $\Rightarrow$ sistema fragile (perde sproporzionalmente con
  lo stress).
\end{itemize}

\Hypoth{} \textbf{Proprietà:} \(A\) è una \emph{proprietà structural}
del sistema, non una risposta post-crisi. È determinata
dall'architettura del sistema (ridondanza, modularità, opzionalità)
first che la crisi si manifesti.

\subsection{\texorpdfstring{Il Coefficiente di Involuzione
\(\Omega\)}{Il Coefficiente di Involuzione \textbackslash Omega}}\label{il-coefficiente-di-involution-omega}

\Proved{} \textbf{Definizione formal:}

\[
\Omega = \frac{\Delta E_{\text{dissipata}}}{\Delta E_{\text{utile}}} - 1
\]

where \(\Delta E\) è il flusso di energia sistemica (risorse, capitale,
attenzione) per unità di tempo. Interpretazione:

\begin{itemize}
\tightlist
\item
  \(\Omega = 0\): ogni unità di energia dissipata produce un'unità di
  valore utile (break-even involutivo);
\item
  \(\Omega > 0\): il sistema dissipa più di quanto produca ---
  \textbf{Neijuan} in senso stretto;
\item
  \(\Omega \to \infty\): collapse energetico-produttivo.
\end{itemize}

\Hypoth{} \textbf{Ancoraggio empirico:} \(\Omega\) è measurable come
complemento della Total Factor Productivity (TFP). Sistemi con TFP in
declino persistente hanno \(\Omega\) crescente.

\Hypoth{} \textbf{Parentela con Turchin:} la \emph{elite overproduction}
misurata come inverso del reddito relativo delle élite è un caso
particolare di \(\Omega\) applicato al sottosistema elitario. Turchin
cattura \(\Omega_{\text{élite}}\); la DEQ$^2$ generalizza.

\Proved{} \textbf{Microfondazione evolutiva: followszione da CPGG.}
Quando la struttura di cattura è descrivibile through il CPGG
(§D.4), \(\Omega\) è \emph{followsbile} dai parameters del gioco ---
non solo definito. All'equilibrio del CPGG, con frazione estrattiva
\(\phi\), quota asimmetrica \(\eta > 1/N\) e moltiplicatore \(r\) del
bene pubblico:

\[
\Omega(\phi,\eta,r) = \frac{\phi\cdot\eta\cdot r}{(1-\phi\eta)r - 1}
\]

Per \(\phi \ll 1\) (estrattori rari ma organizzati):

\[
\boxed{\Omega \approx \frac{\phi\cdot\eta\cdot r}{r-1}} \tag{P4}
\]

Il comportamento limite è rivelatore: per \(r \to 1^+\) (nessun
surplus cooperational), \(\Omega \to \infty\) --- ogni estrazione è
totalmente distruttiva; per \(r \to \infty\), \(\Omega \to \phi\eta\)
(quota asimmetrica pura). \(\Omega\) non è una proprietà dell'attore
singolo ma una funzione del trio (\(\phi\), \(\eta\), \(r\)):
concentrazione estrattiva, asimmetria e produttività collettiva.



\section{\texorpdfstring{Derivazione di
\(T_s^{(2)}\)}{Derivazione di T\_s\^{}\{(2)\}}}\label{followszione-di-t_s2}

\Hypoth{} \textbf{Principio di incorporazione:} \(A\) e \(\Omega\) non
sono variables indipendenti aggiuntive, ma \textbf{modulatori
strutturali} delle quantità già presenti in \(T_s^{(1)}\).

\textbf{Modulazione di \(\sigma\) through \(A\).} Un sistema antifragile
percepisce una volatilità effettiva ridotta, because trasforma parte del
rumore in coerenza:

\[
\sigma_{\text{eff}} = \frac{\sigma}{1 + A}
\]

\begin{itemize}
\tightlist
\item
  \(A > 0\): \(\sigma_{\text{eff}} < \sigma\) (la volatilità è domata);
\item
  \(A = 0\): \(\sigma_{\text{eff}} = \sigma\) (caso neutro, recuperiamo
  la DEQ$^1$);
\item
  \(A \to -1^+\): \(\sigma_{\text{eff}} \to \infty\) (fragility
  catastrofica).
\end{itemize}

\textbf{Modulazione di \(\varepsilon_{\text{dis}}\) through \(\Omega\).}
L'involution amplifica il disimpegno: ogni unità dissipata in
competizione improduttiva aumenta l'eteronomia del sistema
(Bandura/Milgram):

\[
\varepsilon_{\text{eff}} = \varepsilon_{\text{dis}} \cdot (1 + \Omega)
\]

\begin{itemize}
\tightlist
\item
  \(\Omega = 0\): nessuna amplificazione;
\item
  \(\Omega > 0\): disimpegno amplificato in proporzione all'involution.
\end{itemize}

\textbf{Equazione unificata DEQ$^2$ (forma individuale/meso):}

\[
\boxed{T_s^{(2)} = \frac{s_{iii} \cdot z \cdot \delta \cdot (1+A)}{\sigma \cdot \varepsilon_{\text{dis}} \cdot (1+\Omega)}}
\]

\Proved{} \textbf{Verifica di consistenza.} Per \(A = 0\) e
\(\Omega = 0\), \(T_s^{(2)} \to T_s^{(1)}\). La DEQ$^1$ è il caso limite
della DEQ$^2$ in assenza delle due modulazioni.



\section{Proprietà dell'Equazione}\label{proprietuxe0-dellequation}

\subsection{Monotonicità}\label{monotonicituxe0}

\Proved{} Per ogni variable, il segno della followsta è determinato:

\[
\frac{\partial T_s^{(2)}}{\partial s_{iii}},\ \frac{\partial T_s^{(2)}}{\partial z},\ \frac{\partial T_s^{(2)}}{\partial \delta},\ \frac{\partial T_s^{(2)}}{\partial A} > 0
\]

\[
\frac{\partial T_s^{(2)}}{\partial \sigma},\ \frac{\partial T_s^{(2)}}{\partial \varepsilon_{\text{dis}}},\ \frac{\partial T_s^{(2)}}{\partial \Omega} < 0
\]

Le prime sono \textbf{forze di coerenza}, le seconde \textbf{forze di
decoherence}. Nessuna variable ha segno ambiguo --- proprietà importante
per l'interpretabilità.

\subsection{Asintoti Critici}\label{asintoti-critici}

\begin{itemize}
\tightlist
\item
  \(\delta \to 0\): collapse dell'time horizon (miopia totale) $\Rightarrow$
  \(T_s^{(2)} \to 0\).
\item
  \(\sigma \to \infty\): volatilità illimitata $\Rightarrow$ \(T_s^{(2)} \to 0\), a
  meno che \(A \to \infty\) (caso teorico di antifragility infinita).
\item
  \(\Omega \to \infty\): involution totale $\Rightarrow$ \(T_s^{(2)} \to 0\).
\item
  \(A \to -1\): fragility catastrofica $\Rightarrow$ \(T_s^{(2)} \to 0\) anche con
  volatilità finita.
\end{itemize}

\subsection{Punto Critico di Transizione di
Fase}\label{punto-critico-di-transizione-di-fase}

\Hypoth{} \textbf{Ipotesi operativa:} la phase transition del sistema
(collapse, rivoluzione, riorganizzazione structural) si manifesta
quando \(T_s^{(2)} < 1\) per una durata persistente
\(\Delta t > \Delta t^*\), where \(\Delta t^*\) dipende dalla scala del
sistema:

{\def\LTcaptype{none} % do not increment counter
{\small\begin{longtable}[]{@{}ll@{}}
\toprule\noalign{}
Scala & \(\Delta t^*\) stimato \\
\midrule\noalign{}
\endhead
\bottomrule\noalign{}
\endlastfoot
Individuo / azienda & 4-12 mesi \\
Sistema nazionale & 3-7 anni \\
Ciclo egemonico & 15-25 anni (coerente con Arrighi) \\
\end{longtable}}
}

\Conj{} \textbf{Congettura di scala:} \(\Delta t^*\) è
approssimativamente proporzionale alla radice del numero di attori del
sistema --- hypothesis da verificationre empiricamente nei capitoli applicati
(Part III).



\section{Limiti di Questa
Formulazione}\label{limiti-di-questa-formulazione}

\Conj{} \textbf{Limite 1.} \(T_s^{(2)}\) è qui scritto a \emph{livello
local/meso} --- descrive un attore come se le sue componenti fossero
allineate (\(\Phi = 1\)). L'aggregazione formal al sistema globale, già
anticipata nel \S{}3.1bis through l'operator \(\langle\cdot\rangle_\Phi\),
è completata nel Ch.~4 where \(T_s^{(2),\text{sys}}\) assume la sua
forma definitiva con \(\Phi < 1\). La differenza tra \(T_s^{(2)}\) e
\(T_s^{(2),\text{sys}}\) misura quanto un sistema \emph{spreca}
legitimacy per decoherence interna --- il caso paradigmatico è l'USA
contemporaneo.

\Conj{} \textbf{Limite 2.} Le variables
\(s_{iii}, z, \sigma, \varepsilon_{\text{dis}}, A, \Omega\) sono qui
assunte misurabili. La operational definition di \emph{come misurarle} è
rinviata all'\textbf{Appendix B} (Protocolli di measurement dei
parameters).

\ToVerify{} \textbf{Da verificationre:} la scelta della forma moltiplicativa
\((1+A)\) e \((1+\Omega)\) è semplice e monotonicamente consistente, ma
non è l'unica possibile. Forme alternative (esponenziale, razionale)
andrebbero testate su dati storici. Questa verification è rinviata al paper
tecnico in inglese.



\section{Chapter Summary}\label{sintesi-del-chapter}

\Proved{} \textbf{In una frase:} la stabilità di un sistema egemonico è
il rapporto tra una coerenza (legitimacy $\times$ dialettica $\times$ orizzonte)
\emph{amplificata} dall'antifragility e una decoherence (volatilità $\times$
disimpegno) \emph{amplificata} dall'involution.

\[
T_s^{(2)} = \frac{\underbrace{s_{iii} \cdot z \cdot \delta}_{\text{coerenza}} \cdot \overbrace{(1+A)}^{\text{amplificatore}}}{\underbrace{\sigma \cdot \varepsilon_{\text{dis}}}_{\text{decoerenza}} \cdot \underbrace{(1+\Omega)}_{\text{amplificatore}}}
\]

Il Ch.~4 affronterà il passaggio al sistema aggregato introducendo il
order parameter \(\Phi\).



\section{Closing Epistemic Notes
Chapter}\label{note-epistemiche-di-chiusura-chapter}

{\def\LTcaptype{none} % do not increment counter
{\small\begin{longtable}[]{@{}
  >{\raggedright\arraybackslash}p{(\linewidth - 4\tabcolsep) * \real{0.3333}}
  >{\raggedright\arraybackslash}p{(\linewidth - 4\tabcolsep) * \real{0.3333}}
  >{\raggedright\arraybackslash}p{(\linewidth - 4\tabcolsep) * \real{0.3333}}@{}}
\toprule\noalign{}
\begin{minipage}[b]{\linewidth}\raggedright
\#
\end{minipage} & \begin{minipage}[b]{\linewidth}\raggedright
Claim
\end{minipage} & \begin{minipage}[b]{\linewidth}\raggedright
Status
\end{minipage} \\
\midrule\noalign{}
\endhead
\bottomrule\noalign{}
\endlastfoot
1 & La correzione notezionale di \(\delta\) & \Proved{} coerenza
interna \\
2 & L'identità formal tra \(\Pi\) Topografia e \(\Pi\) Geometria
(\(\delta^*=0{,}50\)) & \Proved{} verified sui due testi \\
3 & L'transposition operator \(\langle\cdot\rangle_\Phi\) (Kuramoto)
& \Hypoth{} hypothesis formal di scala \\
4 & La followszione di \(T_s^{(1)}\) dal Memoriale & \Proved{} richiamo
fedele \\
5 & Le definizioni di \(A\) e \(\Omega\) & \Proved{} formalizzazioni
rigorose \\
6 & L'equation \(T_s^{(2)}\) & \Hypoth{} working hypothesis
unificante \\
7 & La condizione \(T_s^{(2)} < 1\) come soglia di transizione &
\Hypoth{} hypothesis testabile \\
8 & La scala \(\Delta t^*\) & \Conj{} conjecture di scala \\
9 & La forma moltiplicativa \((1+A)\), \((1+\Omega)\) & \ToVerify{} da
verificationre empiricamente \\
\end{longtable}}
}


